\section{From completion to point sections}\label{sec:sections}

We now use Theorem~\ref{thm:completion} to compare decomposition groups.
The comparison lifts any section localized at a rank-one valuation to a
birational section.  Koenigsmann's theorem then places the same section at
a rational point, and the two valuations give a contradiction.  The
remaining valuations are horizontal refinements, where the localization
theorem of Pop--Stix already identifies the section.

\begin{corollary}\label{cor:decomposition}
Let $k/\mathbb Q_p$ be finite and let $X/k$ be a smooth proper
geometrically connected curve of genus at least two.  Put $K=k(X)$ and
let $\widetilde K\subset K^{\mathrm{sep}}$ be the extension corresponding
to the full arithmetic \'etale fundamental group of $X$.  Let $w$ be a
real-valued rank-one valuation of $K$ extending the $p$-adic valuation,
normalized by $w(p)=1$, and fix a prolongation $w^{\mathrm{sep}}$ to
$K^{\mathrm{sep}}$.  Restriction induces an isomorphism of profinite
groups
\[
 D_{w^{\mathrm{sep}}}(K^{\mathrm{sep}}/K)
 \xrightarrow{\ \sim\ }
 D_{\widetilde w}(\widetilde K/K),
 \qquad \widetilde w=w^{\mathrm{sep}}|_{\widetilde K}.
\]
\end{corollary}

\begin{proof}
Restriction is surjective.  Indeed, lift an automorphism stabilizing
$\widetilde w$ to an automorphism of $K^{\mathrm{sep}}/K$.  Its translate
of $w^{\mathrm{sep}}$ is another prolongation of $\widetilde w$.
Transitivity of $\operatorname{Gal}(K^{\mathrm{sep}}/\widetilde K)$ on
these prolongations allows us to adjust the lift so that it stabilizes
$w^{\mathrm{sep}}$.  The transitivity statement for an infinite normal
extension follows from its finite-level version by compactness.

An automorphism stabilizing a rank-one valuation preserves its normalized
absolute value.  To see why normalization matters, any two
order-preserving embeddings of the value group into $\mathbb R$ differ
by a positive scale; the value on $p$ fixes that scale.  Such an
automorphism therefore extends as an isometry to
$E=\widehat{K^{\mathrm{sep}}}$.  If its restriction to $\widetilde K$ is
the identity, it fixes $H=\widehat{\widetilde K}\subset E$ pointwise.
Theorem~\ref{thm:completion} gives $K^{\mathrm{sep}}\subset H$, so the
automorphism is the identity.

Both decomposition groups are closed subgroups of their Galois groups.
For example, stabilization of a valuation ring is the intersection,
over its algebraic elements, of conditions on finite Galois orbits, each
of which is clopen.  Restriction is consequently a continuous bijection
between compact Hausdorff groups, and its inverse is continuous.
\end{proof}

\begin{lemma}\label{lem:no-rank-one-section}
With $k$, $X$, $K$, and $\widetilde K$ as in
Corollary~\ref{cor:decomposition}, let $w$ be any rank-one valuation of
$K$ extending the $p$-adic valuation and let $\widetilde w$ be any
prolongation to $\widetilde K$.  The projection
\[
 D_{\widetilde w}(\widetilde K/K)\longrightarrow G_k
\]
has no continuous group-homomorphic section.
\end{lemma}

\begin{proof}
Normalize $w(p)=1$ and prolong the valuation to $K^{\mathrm{sep}}$.
Suppose that a section exists.  Corollary~\ref{cor:decomposition} lifts
it to a continuous section
\[
 \widehat s:G_k\longrightarrow G_K
       =\operatorname{Gal}(K^{\mathrm{sep}}/K)
\]
whose image stabilizes this prolongation.  The birational section
theorem of Koenigsmann places the entire image of $\widehat s$ also in
the decomposition group of a prolongation of $\operatorname{ord}_a$ for
some $a\in X(k)$ \cite[preprint version, Proposition~2.4(2), Definition~1.4, and
Observation~1.3]{Koenigsmann}.  Here $\operatorname{ord}_a$ is the
closed-point valuation of $K$, trivial on $k$.

Write $\Sigma=\widehat s(G_k)$ and $M=(K^{\mathrm{sep}})^\Sigma$.
The subgroup $\Sigma$ is compact and hence closed, so
$G_M=\Sigma$.  Let $v_1$ and $v_2$ be the restrictions to $M$ of the two
selected prolongations.  Each is henselian: its absolute Galois group
stabilizes a prolongation, and the transitive action on all
prolongations therefore has a single point.  Both valuations are
nontrivial and have rank one.  Indeed, value groups in an algebraic
valued extension differ by a torsion quotient, which preserves their
convex-subgroup chains.  Moreover,
\[
 v_1(p)=1,\qquad v_2|_{k^\times}=0,
\]
so these valuations are inequivalent.  The resulting contradiction is
an instance of the classical uniqueness principle for henselian
valuations associated with F.~K.~Schmidt \cite{SchmidtValuations}.
We give explicitly the two-valuation approximation needed here.

Choose a nonsquare $d\in k^\times$.  It remains a nonsquare in $M$:
otherwise $G_M=\Sigma$ would fix a square root of $d$ in $\bar k$,
contrary to the surjection $\Sigma\to G_k$.  Let $t\in K$ be a
uniformizer at $a$.  Choose a positive integer $N$ with $N>v_1(t)$,
and put $z=p^N/t$.  Then
\[
 v_1(z)>0,\qquad v_2(z)=-v_2(t)<0.
\]
For positive integers $r$, set
\[
 q_r=d+(1-d)\frac{z^r}{1+z^r}.
\]
The denominator is nonzero because $v_1(z^r)>0$.  At the first
valuation $q_r$ tends to $d$, and at the second it tends to $1$.
More precisely,
\[
 v_1(q_r/d-1)
   =v_1((1-d)/d)+r v_1(z)\longrightarrow+\infty,
\]
whereas
\[
 v_2(q_r-1)
   =v_2(d-1)-r v_2(z)\longrightarrow+\infty.
\]
For sufficiently large $r$, Hensel's lemma applied at $1$ gives a
square root of $q_r/d$ using $v_1$, and a square root of $q_r$ using
$v_2$.  The sufficient inequalities are
\[
 v_1(q_r/d-1)>2v_1(2),\qquad
 v_2(q_r-1)>2v_2(2).
\]
Their quotient makes $d$ a square in $M$, a contradiction.  This
argument includes $p=2$ and arbitrary nondiscrete rank-one value groups.
\end{proof}

\begin{proof}[Proof of the surjectivity in Theorem~\ref{thm:section}]
Let $s:G_k\to\Pi_X$ be a continuous group-homomorphic section.
By the valuative localization theorem of Pop--Stix
\cite[preprint version, Theorem~26]{PopStix}, there is a valuation $w$ of $K$ extending
the $p$-adic valuation, together with a prolongation $\widetilde w$ to
$\widetilde K$, such that
\[
 s(G_k)\subseteq D_{\widetilde w}(\widetilde K/K).
\]
We classify the possibilities for $w$.

The valuation transcendence inequality gives
\[
 \operatorname{ratrank}(\Gamma_w/\Gamma_k)\leq 1.
\]
Since $\Gamma_k$ has rational rank one and ordered rank is at most
rational rank, $w$ has rank at most two.  Rank one is excluded by
Lemma~\ref{lem:no-rank-one-section}.  If $w$ has rank two, let $u$ be
its rank-one coarsening.  The nonzero proper convex subgroup of the
rank-two value group has a unique extension to the value group of
$\widetilde w$: it consists of those values some positive multiple of
which lies in the original subgroup.  This follows because an algebraic
valued extension has torsion quotient of value groups.  Every
automorphism stabilizing $\widetilde w$ preserves this subgroup, so
$s(G_k)$ stabilizes the resulting prolongation of $u$.  If $u$
is nontrivial on $k$, it can be normalized to extend the $p$-adic
valuation and is excluded by the same lemma.

It remains that $u$ is trivial on $k$.  Properness gives it a center
on $X$.  This center cannot be generic because $u$ is nontrivial, so
it is a closed point $a$.  The valuation ring of $u$ dominates the
discrete valuation ring $\mathcal O_{X,a}$.  Every element of $K^\times$
is a power of a uniformizer times a unit of this local ring; hence
$u$ is equivalent to $\operatorname{ord}_a$.  Its residue field is the
finite extension $k(a)/k$.  The second valuation in the composition
$w$ is the unique prolongation of the $p$-adic valuation to $k(a)$.
Thus $w$ is of type $2h$.

The chosen horizontal prolongation to $K\bar k$ selects a geometric
point above $a$.  Its stabilizer acts on constants through
$G_{k(a)}\subseteq G_k$, with the embedding determined by that point.
Since $s(G_k)$ projects onto $G_k$, necessarily $k(a)=k$.
For the full \'etale tower the horizontal inertia is trivial, and the
horizontal decomposition group is the point section at $a$.
Equivalently, the type-$2h$ conclusion of
\cite[preprint version, Section~6.1]{PopStix} gives $s=s_a$ after the corresponding
geometric choice of base point.  This proves surjectivity.

Notice that a type-$2h$ valuation can lie over a non-rational closed
point.  Rationality here follows from the surjection to $G_k$; no
unrestricted equivalence between ``not type $2h$'' and ``not originating
from a rational point'' is used.
\end{proof}

\begin{proof}[Proof of the injectivity in Theorem~\ref{thm:section}]
Suppose $a,b\in X(k)$ have the same section class modulo the full
geometric fundamental group.  Let $J$ be the Jacobian of $X$, and use
$a$ as base point for the Abel--Jacobi map
\[
 i_a:X\longrightarrow J,\qquad x\longmapsto [x-a].
\]
Their induced sections on $J$ are again geometrically conjugate.
For every positive integer $n$, multiplication
$[n]:J\to J$ is a finite \'etale torsor under $J[n]$, also when
$p\mid n$ since $k$ has characteristic zero.  Conjugate sections give
isomorphic $G_k$-actions on its geometric fiber, so the existence of a
fixed point is unchanged by conjugation.  The section at $i_a(a)=0$
has the fixed lift $0$.  Therefore the fiber over $i_a(b)$ has a
$G_k$-fixed point, which is a $k$-rational point of that fiber.  Thus
\[
 i_a(b)\in\bigcap_{n\geq1}nJ(k).
\]
This is also the usual Kummer argument: the class of the multiplication
torsor pulled back along a section is unchanged by geometric conjugation.

The compact $p$-adic analytic group $J(k)$ has an open subgroup
$U\cong\mathbb Z_p^{g(X)[k:\mathbb Q_p]}$ by the local logarithm.  If
$q=[J(k):U]$, then $qJ(k)\subseteq U$.  For each $r\geq0$, take
$n=qp^r$ in the displayed intersection.  Any element of that intersection
can then be written as $p^r(qy)$ with $y\in J(k)$, and hence belongs
to $p^rU$.  As $U$ is a finite free $\mathbb Z_p$-module,
\[
 \bigcap_{r\geq0}p^rU=0,
\]
and hence $i_a(b)=0$.  If $a\ne b$, the linear equivalence of the
degree-one divisors $a$ and $b$ would give a degree-one morphism
$X\to\mathbb P^1$, forcing genus zero.  Hence $a=b$, as required.
\end{proof}
